%%=============================================================================
%% Voorwoord
%%=============================================================================

\chapter*{\IfLanguageName{dutch}{Woord vooraf}{Preface}}%
\label{ch:voorwoord}

%% Het voorwoord is het enige deel van de bachelorproef waar je vanuit je
%% eigen standpunt (``ik-vorm'') mag schrijven. Je kan hier bv. motiveren
%% waarom jij het onderwerp wil bespreken.
%% Vergeet ook niet te bedanken wie je geholpen/gesteund/... heeft

Het idee voor deze bachelorproef is ontstaan door een combinatie van verschillende zaken. Mijn interesse in automatisatie, en de grote focus in de mainframe industrie op modernisatie en het tekort aan werkkrachten, vormden de basis van het idee. Daarnaast had ik al eerdere ervaring met Rexx en Python.

Eerst en vooral wil ik graag Leendert Blondeel bedanken. Hij is de spil van de mainframe opleiding aan de Hogent. Hij introduceerde me in het mainframe landschap en daar ben ik hem erg dankbaar voor. 

Nadat ik had besloten om een vergelijking van Rexx en Python te doen, heb ik in de mainframe community op Discord gevraagd voor advies en ideeën. Ik wil graag Arthur Coucke bedanken. Hij kwam met het idee om het genereren van JCL te gebruiken in de vergelijking. Na enkele e-mails stelde Arthur voor om, in plaats van enkel Rexx en Python, het hele automatisatie landschap te analyseren. Uiteindelijk heb ik, na een gesprek met mijn begeleider, Leendert Blondeel, besloten om de vergelijking enkel tussen Rexx en Python te houden, maar in plaats van één grote PoC, verschillende kleinere PoC's te combineren om zo een uitgebreidere analyse te kunnen doen van Rexx en Python. Eén van die PoC's is het genereren van JCL.

Ik zou graag mijn co-promotor Bobby Tjassens Keiser bedanken. Bij het uitwerken van mijn thesis heb ik altijd kunnen rekenen op zijn feedback en expertise. Voor feedback op de inhoud, maar ook voor code voorbeelden kon ik doorheen heel het proces op de ervaring van Bobby rekenen. Via wekelijkse meetings, en toegang tot een systeem om op te werken, heeft Bobby deze bachelorproef voor een groot deel mogelijk gemaakt. Ik heb enorm veel bijgeleerd over automatisatie, maar ook over mainframes in het algemeen. Deze bachelorproef was niet mogelijk geweest zonder de hulp van Bobby.

Een groep mensen die ik ook zeker niet mag vergeten bedanken, is de mainframe community, zowel op Discord als op Linkedin. Ik heb enorm veel input en feedback gekregen op elke vraag die ik heb gesteld. Ook het aantal antwoorden, en de kwaliteit ervan, op mijn enquête heeft me enorm vooruit geholpen om de beste requirements en use cases te bepalen. Ook wil ik graag Emma Skovgård bedanken voor haar tijd en input tijdens het interview, dat mede aan de basis lag van mijn PoC over RACF.

Tot slot wil ik ook mijn zus bedanken, die de structuur en taal van mijn bachelorproef enorm veel heeft vooruitgeholpen en verbeterd. Voor het schrijven van deze bachelorproef is geen artificiële intelligentie gebruikt, in plaats daarvan had ik mijn zus.